\subsection*{الف}
چون همهٔ هدها projectionهای یکسان دارند، تمام هدها خروجی یکسانی تولید می‌کنند:
\begin{flalign*}
    W_i^Q = W^Q,\; W_i^K = W^K,\; W_i^V = W^V \;\;\forall i &&\\
    \implies \text{head}_i = \text{Attention}(QW^Q, KW^K, VW^V) =: H \quad \forall i &&
\end{flalign*}
پس هر $h$ هد دقیقاً همان attention را $h$ بار محاسبه می‌کنند. خروجی لایه:
\begin{flalign*}
    \text{MHA}(Q,K,V) &= \text{Concat}(H, H, \dots, H)\, W^O &&\\
    &= \underbrace{[\,H\,|\,H\,|\,\cdots\,|\,H\,]}_{h \text{ copies}}\, W^O &&
\end{flalign*}
با افراز سطری $W^O \in \mathbb{R}^{(h\,d_k)\times d_{\text{model}}}$ به $h$ بلوک $W^O = \begin{bmatrix} W^O_{(1)} \\ \vdots \\ W^O_{(h)} \end{bmatrix}$، خاصیت Concat–ضرب بلوکی می‌دهد:
\begin{flalign*}
    \text{MHA}(Q,K,V) &= \sum_{i=1}^{h} H\, W^O_{(i)} = H\Big(\sum_{i=1}^{h} W^O_{(i)}\Big) = H\, \widetilde{W}^O &&
\end{flalign*}
که در آن $\widetilde{W}^O := \sum_{i=1}^h W^O_{(i)} \in \mathbb{R}^{d_k \times d_{\text{model}}}$. پس خروجی همان یک single-head attention است که با یک projection مؤثر $\widetilde{W}^O$ ترکیب شده.

توجه کنید این لزوماً مضرب اسکالر $H$ نیست؛ $\widetilde{W}^O$ یک ماتریس دلخواه است و می‌تواند ساختار غیربدیهی داشته باشد. تنها در حالت خاصِ $\widetilde{W}^O = c\,I$ به مقیاس ساده تقلیل می‌یابد.

نتیجه: قدرت بیانی Multi-Head نه از تعدد هدها، بلکه از \lr{diversity} بین آن‌ها می‌آید. هدهای یکسان \lr{rank} ترکیب را افزایش نمی‌دهند و کل لایه تا یک projection خطی به یک هد فرومی‌پاشد؛ پس تنوع subspaceها شرط لازمِ بهره‌بردن از چندهد بودن است.

\subsection*{ب}
معیار اهمیت تک‌هدی را از طریق حساسیت task loss به حذف هد تعریف می‌کنیم (به سبک \lr{Michel et al. 2019})، با gate $\xi_i \in \{0,1\}$ روی هد $i$:
\begin{flalign*}
    I_i &= \mathbb{E}_{x}\left| \mathcal{L}(x \mid \xi_i = 1) - \mathcal{L}(x \mid \xi_i = 0) \right| &&\\
    &\approx \mathbb{E}_{x}\left| \frac{\partial \mathcal{L}}{\partial \xi_i} \right|
    = \mathbb{E}_{x}\left| H_i^{\top} \frac{\partial \mathcal{L}}{\partial H_i} \right| &&
\end{flalign*}
که از بسط مرتبهٔ اول $\mathcal{L}$ حول $\xi_i = 1$ به دست می‌آید و در یک پاس \lr{backward} برای همهٔ هدها قابل محاسبه است.

برای جنبهٔ افزونگی، شباهت زوجی هدها را با $R(i,j) = \cos(\text{vec}(A_i), \text{vec}(A_j))$ می‌سنجیم. هدی که هم اهمیت پایین دارد و هم با هدی دیگر به‌شدت همبسته است، نامزد حذف است. معیار ترکیبی اولویتِ حذف:
\begin{flalign*}
    \text{score}(i) = \underbrace{I_i}_{\text{اهمیت}} - \lambda \underbrace{\max_{j \neq i,\, j \notin \text{pruned}} R(i,j)}_{\text{افزونگی}} &&
\end{flalign*}
هدها را بر حسب $\text{score}$ صعودی مرتب کرده و $h-k$ هد با کمترین مقدار را حذف می‌کنیم. عبارت دوم تضمین می‌کند از هر خوشهٔ هدهای مشابه فقط یک نماینده باقی بماند (نه آنکه چند نسخهٔ مهمِ مشابه همگی نگه‌داشته شوند).

\subsection*{ج}
متغیر دروازه‌ای $z_i \in [0,1]$ را با تابع هزینه‌ای سه‌بخشی آموزش می‌دهیم:
\begin{flalign*}
    \mathcal{J} = \underbrace{\mathcal{L}_{\text{task}}(z)}_{\text{(۱) عملکرد}}
    \;+\; \alpha \underbrace{\sum_{i=1}^{h} z_i}_{\text{(۲) تنکی}}
    \;+\; \beta \underbrace{\sum_{i<j} z_i\, z_j\, R(i,j)}_{\text{(۳) افزونگی}} &&
\end{flalign*}
نقش هر جمله:
\begin{flalign*}
    &\bullet\;\; \mathcal{L}_{\text{task}}(z): \text{ افتِ عملکرد را جریمه می‌کند} \Rightarrow z_i \text{ روی هدهای مفید را بالا نگه می‌دارد.} &&\\
    &\bullet\;\; \alpha \textstyle\sum_i z_i \;(\ell_1): \text{ فشارِ تنکی} \Rightarrow z_i \text{ها را به سمت صفر می‌راند (حذف هدهای کم‌اهمیت).} &&\\
    &\bullet\;\; \beta \textstyle\sum_{i<j} z_i z_j R(i,j): \text{ جریمهٔ زوجی} \Rightarrow \text{نگه‌داشتن همزمان دو هدِ مشابه را گران می‌کند.} &&
\end{flalign*}
از آنجا که هدف بردارِ تنک می‌خواهد، $\ell_1$ مناسب‌تر از $\ell_2$ است؛ و برای آنکه gate واقعاً به صفر/یک نزدیک شود می‌توان $z_i$ را با \lr{Hard-Concrete} پارامتری کرد تا $\sum_i z_i$ تخمینی از $\|z\|_0$ شود.

انتخاب $k$ پس از آموزش: هدها را بر حسب $z_i$ نزولی مرتب می‌کنیم. مقادیر $z_i$ معمولاً قطبیده می‌شوند (نزدیک ۰ یا ۱)؛ شکافِ طبیعی در طیف مرتب‌شدهٔ $z_i$ مرز را مشخص می‌کند و $k$ را برابر تعداد هدهای بالای آستانه (مثلاً $z_i > 0.5$) می‌گیریم — یعنی $k$ به‌جای آنکه از پیش ثابت شود، از خودِ داده استخراج می‌گردد.

تحلیل دو هدِ کاملاً یکسان ($R(i,j) = 1$، رفتار همسان): جملهٔ افزونگی روی این زوج برابر $\beta\, z_i z_j$ است که با نگه‌داشتن همزمانِ هر دو ($z_i = z_j = 1$) بیشینه و برابر $\beta$ می‌شود. اما چون رفتارشان یکسان است، خروجی $z_i H_i + z_j H_j = (z_i + z_j) H$ تنها به مجموع $z_i + z_j$ وابسته است؛ پس همان عملکرد با $(z_i, z_j) = (1, 0)$ نیز حاصل می‌شود بی‌آنکه $\mathcal{L}_{\text{task}}$ تغییر کند، حال‌آنکه هم جریمهٔ تنکی ($-\alpha$) و هم جریمهٔ افزونگی ($-\beta$) کاهش می‌یابند:
\begin{flalign*}
    \mathcal{J}(1,1) - \mathcal{J}(1,0) = \alpha + \beta > 0 &&
\end{flalign*}
بنابراین جریمهٔ افزونگی $\mathcal{J}$ را در حالتِ «نگه‌داشتن هر دو» اکیداً بزرگ‌تر می‌کند بی‌آنکه نفعی در task بدهد؛ پس راه‌حل بهینه یکی از این دو هدِ تکراری را خاموش می‌کند و نگه‌داشتن همزمان هر دو نامطلوب است.